\documentclass[a4paper, twoside]{report}

%% Language and font encodings
\usepackage[english]{babel}
\usepackage[utf8x]{inputenc}
\usepackage[T1]{fontenc}
\usepackage{mathrsfs}

%% Sets page size and margins
\usepackage[a4paper,top=3cm,bottom=2cm,left=3cm,right=3cm,marginparwidth=1.75cm]{geometry}

\usepackage{amsthm}
\usepackage{graphicx}
\usepackage[]{listings}
\usepackage{booktabs}
\usepackage[colorinlistoftodos]{todonotes}
\usepackage[colorlinks=true, allcolors=blue]{hyperref}
\usepackage{comment}
\usepackage[numbers]{natbib}
\usepackage{fancyref}
\usepackage{wrapfig}

%% Math
\usepackage{mathtools}
\usepackage{stmaryrd}

\usepackage{title/titlestructure}

%% Utils
\usepackage{background/utils}


\title{JOINBENCH: Designing tools for the evaluation of efficient equijoins in Haskell}
\author{Matteo Bongiovanni}
\date{January 2023}
\supervisor{Dr.\ Nicolas Wu}
\secondmarker{Dr.\ Steffen van Bakel}
\documenttype{MEng Individual Project}
\university{Imperial College London}
\department{Department of Computing}
\degree{Joint Mathematics and Computing}

\begin{document}
\begin{comment}
General Information
The project report (occasionally referred to as a thesis) is an extremely
important aspect of the project. It serves to show what you have achieved and
should demonstrate that:

You understand the wider context of computing by relating your choice of project,
and the approach you take, to existing products or research.
You can apply the theoretical and practical techniques taught in the course to the
problem you are addressing and that you understand their relevance to the wider
world of computing.
You are capable of objectively criticising your own work and making constructive
suggestions for improvements or further work based on your experiences so far.
As a computing professional, you operate ethically and can explain your thinking
and working processes clearly and concisely to third parties who may not be
experts in the field in which you are working.
Most of the project assessors will not have followed the project throughout and
will only have a short time to listen to a presentation or see a demonstration.
For this reason they will rely heavily on the report to judge the project

Many students underestimate the importance of the report and make the mistake of
thinking that top marks can be achieved simply for producing a good product. This
is fundamentally not the case and many projects have been graded well below their
potential because of an indifferent or poor write-up. In order to get the balance
right you should consider that the aim of the project is to produce a good report
and that software, hardware, theory etc. that you developed during the project are
merely a means to this end. Don't make the mistake of leaving the write-up to the
last minute. Ideally you should produce the bulk of the report as you go along and
use the last week or two to bring it together into a coherent document.

It is helpful to get feedback from your supervisor about your project report, but
supervisors cannot be expected to look at documents written at the last minute or
at more than one chapter at a time. Allow plenty of time for this

Project Report Length
The nominal maximum report length is 60 pages, plus appendices. Shorter reports
may be appropriate for many projects and longer reports can be submitted with
prior approval (see below).

The page count includes:

title page
abstract
dedications
table of contents
The page count does not include:

bibliography
appendices (but the report should make sense without the appendices)
When writing up your work be reminded of the fact that a 60-page report isn't
necessarily better than a 30-page report, in the same way that a 10,000 line
implementation isn't necessarily better 5,000 line one. Conciseness, clarity and
elegance are invaluable qualities in report writing, just as they are in
programming. Indeed, you will lose credit if it is evident to the assessors that
the report has been padded with superfluous content in order to make it appear
more substantial than the work justifies. The assessors greatly value quality over
quantity.

It is important to understand that the report is assessed on the quality of the
technical writing. Here this means your ability to document your ideas, design,
implementation, evaluation, conclusions etc. succinctly, using appropriate
scientific language, as might a researcher when writing up their findings in an
technical magazine, academic journal or conference proceedings.

If you have good reason to use more than 60 pages you should consult your
supervisor and second marker and obtain their approval in advance - they will be
required to read and assess the whole report, whereas other assessment team
members may only scan key sections.

Project Layout
The physical layout and formatting of the report is also important. A tidy, well
laid out and consistently formatted document makes for easier reading and is
suggestive of a careful and professional attitude towards its preparation. Both
LaTeX and MS Word will allow you to produce a cleanly formatted document. Many
supervisors will advise you to use LaTeX as this solves most of the formatting
problems for you. If your report is to contain a substantial number of
mathematical formulae you are strongly advised to use LaTeX.

Project Report Structure
The exact structure of your project (the chapter titles etc) is up to you. Discuss
with your supervisor what structure works best for your project

A project report will usually contain a number of chapters (Chapter 1, Chapter 2
etc). Each chapter contains Sections (2.1, 2.2 etc) and may also contain
subsections (2.1.1, 2.1.2) etc. Try to avoid too many levels of subheading - three
is usually sufficient.

Projects will usually contain all of the following elements, although these may
not exactly map to chapters.
\end{comment}

% Front matter
\begin{titlepage}

\newcommand{\HRule}{\rule{\linewidth}{0.5mm}} % Defines a new command for the horizontal lines, change thickness here

%----------------------------------------------------------------------------------------
%	LOGO SECTION
%----------------------------------------------------------------------------------------

\includegraphics[width=8cm, draft=false]{title/logo.eps}\\[1cm] % Include a department/university logo - this will require the graphicx package
 
%----------------------------------------------------------------------------------------

\center\ % Center everything on the page

%----------------------------------------------------------------------------------------
%	HEADING SECTIONS
%----------------------------------------------------------------------------------------

\makeatletter
\textsc{\LARGE \@documenttype}\\[1.5cm]
\textsc{\Large \@university}\\[0.5cm]
\ifdefined\@degree\
\textsc{\large \@degree}\\[0.5cm]
\fi
\textsc{\large \@department}\\[0.5cm]
\makeatother

%----------------------------------------------------------------------------------------
%	TITLE SECTION
%----------------------------------------------------------------------------------------
\makeatletter
\HRule\\[0.4cm]
{ \huge \bfseries \@title}\\[0.4cm] % Title of your document
\HRule\\[1.5cm]
 
%----------------------------------------------------------------------------------------
%	AUTHOR SECTION
%----------------------------------------------------------------------------------------

\begin{minipage}{0.4\textwidth}
\begin{flushleft} \large
\emph{Author:}\\
\@author\ % Your name
\end{flushleft}
\end{minipage}
~\begin{minipage}{0.4\textwidth}
\begin{flushright} \large
\emph{Supervisor:} \\
\@supervisor\\[1.2em]

\ifdefined\@secondmarker\
\emph{Second Marker:} \\
\@secondmarker\
\fi

\end{flushright}
\end{minipage}\\[2cm]

%----------------------------------------------------------------------------------------
%	DATE SECTION
%----------------------------------------------------------------------------------------

{\large \@date}\\[2cm] % Date, change the \today to a set date if you want to be precise

\vfill % Fill the rest of the page with whitespace
\makeatother

\end{titlepage}
\begin{comment}
The abstract is a very brief summary of the report's contents. It should be no more than half a page long. Somebody unfamiliar with your project should have a good idea of what it's about having read the abstract alone and will know whether it will be of interest to them. Note that the abstract is a summary of the entire project including its conclusions. A common mistake is to provide only introductory elements in the abstract without saying what has been achieved.
\end{comment}


\begin{abstract}
In order to evaluate alternative equijoin implementation in Haskell, this paper
introduced the \relation{JOINBENCH} relation and surrounding methodology and
tooling. The \relation{JOINBENCH} relation is a scheme designed for the
synthesis of data optimised to help evaluate the performance of equijoins in a
variety of scenarios. Furthermore, this paper presents a low-level library that
helps users customise and define their own synthetic data sources for future
benchmarking purposes.

In their distinguished paper
\relalg{}~\cite{RelationalAlgebraByWayOfAdjunctions} it was noted that the
monadic structure of bulk types can help explain most of relational algebra.
Using this structure, the authors designed a new method to facilitate
the use of monad comprehensions in an efficient implementation of equijoins of
relational databases. This project presents an implementation of such a system
and an evaluation of the performance such query optimisations carry using the
bespoke tooling described above. 
\end{abstract}

\begin{comment}
Table of contents:
This should list the main chapters, sections and subsections of your report.
Choose self-explanatory chapter and section titles and use double spacing for
clarity. If possible you should include page numbers indicating where each chapter
section begins. 
\end{comment}
\documentclass[a4paper, twoside]{report}

%% Language and font encodings
\usepackage[english]{babel}
\usepackage[utf8x]{inputenc}
\usepackage[T1]{fontenc}
\usepackage{mathrsfs}

%% Sets page size and margins
\usepackage[a4paper,top=3cm,bottom=2cm,left=3cm,right=3cm,marginparwidth=1.75cm]{geometry}

\usepackage{amsthm}
\usepackage{graphicx}
\usepackage[]{listings}
\usepackage{booktabs}
\usepackage[colorinlistoftodos]{todonotes}
\usepackage[colorlinks=true, allcolors=blue]{hyperref}
\usepackage{comment}
\usepackage[numbers]{natbib}
\usepackage{fancyref}
\usepackage{wrapfig}

%% Math
\usepackage{mathtools}
\usepackage{stmaryrd}

\usepackage{title/titlestructure}

%% Utils
\usepackage{background/utils}


\title{A Deeper Dive into Relational Algebra by Way of Adjunctions}
\author{Matteo Bongiovanni}
\date{January 2023}
\supervisor{Dr.\ Nicolas Wu}
\secondmarker{Dr.\ Steffen van Bakel}
\documenttype{MEng Individual Project}
\university{Imperial College London}
\department{Department of Computing}
\degree{Joint Mathematics and Computing}

\begin{document}
\begin{comment}
General Information
The project report (occasionally referred to as a thesis) is an extremely important aspect of the project. It serves to show what you have achieved and should demonstrate that:

You understand the wider context of computing by relating your choice of project, and the approach you take, to existing products or research.
You can apply the theoretical and practical techniques taught in the course to the problem you are addressing and that you understand their relevance to the wider world of computing.
You are capable of objectively criticising your own work and making constructive suggestions for improvements or further work based on your experiences so far.
As a computing professional, you operate ethically and can explain your thinking and working processes clearly and concisely to third parties who may not be experts in the field in which you are working.
Most of the project assessors will not have followed the project throughout and will only have a short time to listen to a presentation or see a demonstration. For this reason they will rely heavily on the report to judge the project.

Many students underestimate the importance of the report and make the mistake of thinking that top marks can be achieved simply for producing a good product. This is fundamentally not the case and many projects have been graded well below their potential because of an indifferent or poor write-up. In order to get the balance right you should consider that the aim of the project is to produce a good report and that software, hardware, theory etc. that you developed during the project are merely a means to this end. Don't make the mistake of leaving the write-up to the last minute. Ideally you should produce the bulk of the report as you go along and use the last week or two to bring it together into a coherent document.

It is helpful to get feedback from your supervisor about your project report, but supervisors cannot be expected to look at documents written at the last minute or at more than one chapter at a time. Allow plenty of time for this.

Project Report Length
The nominal maximum report length is 60 pages, plus appendices. Shorter reports may be appropriate for many projects and longer reports can be submitted with prior approval (see below).

The page count includes:

title page
abstract
dedications
table of contents
The page count does not include:

bibliography
appendices (but the report should make sense without the appendices)
When writing up your work be reminded of the fact that a 60-page report isn't necessarily better than a 30-page report, in the same way that a 10,000 line implementation isn't necessarily better 5,000 line one. Conciseness, clarity and elegance are invaluable qualities in report writing, just as they are in programming. Indeed, you will lose credit if it is evident to the assessors that the report has been padded with superfluous content in order to make it appear more substantial than the work justifies. The assessors greatly value quality over quantity.

It is important to understand that the report is assessed on the quality of the technical writing. Here this means your ability to document your ideas, design, implementation, evaluation, conclusions etc. succinctly, using appropriate scientific language, as might a researcher when writing up their findings in an technical magazine, academic journal or conference proceedings.

If you have good reason to use more than 60 pages you should consult your supervisor and second marker and obtain their approval in advance - they will be required to read and assess the whole report, whereas other assessment team members may only scan key sections.

Project Layout
The physical layout and formatting of the report is also important. A tidy, well laid out and consistently formatted document makes for easier reading and is suggestive of a careful and professional attitude towards its preparation. Both LaTeX and MS Word will allow you to produce a cleanly formatted document. Many supervisors will advise you to use LaTeX as this solves most of the formatting problems for you. If your report is to contain a substantial number of mathematical formulae you are strongly advised to use LaTeX.

Project Report Structure
The exact structure of your project (the chapter titles etc) is up to you. Discuss with your supervisor what structure works best for your project.

A project report will usually contain a number of chapters (Chapter 1, Chapter 2 etc). Each chapter contains Sections (2.1, 2.2 etc) and may also contain subsections (2.1.1, 2.1.2) etc. Try to avoid too many levels of subheading - three is usually sufficient.

Projects will usually contain all of the following elements, although these may not exactly map to chapters.
\end{comment}

% Front matter
\begin{comment}
This should include the project title and the name of the author of the report. You can also list the name of your supervisor if you wish. 
\end{comment}

\begin{titlepage}

\newcommand{\HRule}{\rule{\linewidth}{0.5mm}} % Defines a new command for the horizontal lines, change thickness here

%----------------------------------------------------------------------------------------
%	LOGO SECTION
%----------------------------------------------------------------------------------------

\includegraphics[width=8cm, draft=false]{report/title/logo.eps}\\[1cm] % Include a department/university logo - this will require the graphicx package
 
%----------------------------------------------------------------------------------------

\center\ % Center everything on the page

%----------------------------------------------------------------------------------------
%	HEADING SECTIONS
%----------------------------------------------------------------------------------------

\makeatletter
\textsc{\LARGE \@documenttype}\\[1.5cm]
\textsc{\Large \@university}\\[0.5cm]
\ifdefined\@degree\
\textsc{\large \@degree}\\[0.5cm]
\fi
\textsc{\large \@department}\\[0.5cm]
\makeatother

%----------------------------------------------------------------------------------------
%	TITLE SECTION
%----------------------------------------------------------------------------------------
\makeatletter
\HRule\\[0.4cm]
{ \huge \bfseries \@title}\\[0.4cm] % Title of your document
\HRule\\[1.5cm]
 
%----------------------------------------------------------------------------------------
%	AUTHOR SECTION
%----------------------------------------------------------------------------------------

\begin{minipage}{0.4\textwidth}
\begin{flushleft} \large
\emph{Author:}\\
\@author\ % Your name
\end{flushleft}
\end{minipage}
~\begin{minipage}{0.4\textwidth}
\begin{flushright} \large
\emph{Supervisor:} \\
\@supervisor\\[1.2em]

\ifdefined\@secondmarker\
\emph{Second Marker:} \\
\@secondmarker\
\fi

\end{flushright}
\end{minipage}\\[2cm]

%----------------------------------------------------------------------------------------
%	DATE SECTION
%----------------------------------------------------------------------------------------

{\large \@date}\\[2cm] % Date, change the \today to a set date if you want to be precise

\vfill % Fill the rest of the page with whitespace
\makeatother

\end{titlepage}
\begin{comment}
The abstract is a very brief summary of the report's contents. It should be no more than half a page long. Somebody unfamiliar with your project should have a good idea of what it's about having read the abstract alone and will know whether it will be of interest to them. Note that the abstract is a summary of the entire project including its conclusions. A common mistake is to provide only introductory elements in the abstract without saying what has been achieved.
\end{comment}


\begin{abstract}
In order to evaluate alternative equijoin implementation in Haskell, this paper
introduced the \relation{JOINBENCH} relation and surrounding methodology and
tooling. The \relation{JOINBENCH} relation is a scheme designed for the
synthesis of data optimised to help evaluate the performance of equijoins in a
variety of scenarios. Furthermore, this paper presents a low-level library that
helps users customise and define their own synthetic data sources for future
benchmarking purposes.

In their distinguished paper
\relalg{}~\cite{RelationalAlgebraByWayOfAdjunctions} it was noted that the
monadic structure of bulk types can help explain most of relational algebra.
Using this structure, the authors designed a new method to facilitate
the use of monad comprehensions in an efficient implementation of equijoins of
relational databases. This project presents an implementation of such a system
and an evaluation of the performance such query optimisations carry using the
bespoke tooling described above. 
\end{abstract}

\tableofcontents

% Main matter 
\chapter{Introduction} % 1-3 pages
\begin{comment}
It’s a good idea to *try* to write the introduction to your final report early on in the project. However, you will find it hard, as you won’t yet have a complete story and you won’t know what your main contributions are going to be. However, the exercise is useful as it will tell you what you *don’t* yet know and thus what questions your project should aim to answer. For the interim report this section should be a short, succinct, summary of the project’s main objectives. Some of this material may be re-usable in your final report, but the chances are that your final introduction will be quite different.  You are therefore advised to keep this part of the interim report short, focusing on the following questions: What is the problem, why is it interesting and what’s your main idea for solving it?  (DON'T use those three questions as subheadings however!  The answers should emerge from what you write.)
\end{comment}

Databases are absolutely vital to modern day society and contain domain specific knowledge about anything from specialised images of eyes \todo{add citation here} to the structure of crystals \cite{CambridgeStructuralDatabase}. Since its conception many different data models describing how to hold the data in databases have emerged, including the relational model \cite{RelationalModel} and the semi-structured model\cite{DatabaseSystems}.

In this project we concern ourselves with the relational model (to be introduced in \fref{sec:relationalmodel}) as a way of modelling the database. This rich model has many methods of expressing queries, especially relational algebra and relational calculus, both with their strengths and weaknesses \cite{RelationalCalculus,RelationalModel}. \todo{they are equivalent though?} However, the favoured specification of the authors of \cite{RelationalAlgebraByWayOfAdjunctions} seemed to be through list comprehensions, a beautiful feature that ``provide for a concise and expressive notation for writing list-processing code''. \cite{MonadComprehensions} They eventually propose using GHC's extended list comprehension syntax specifically designed to help bridge the already close relationship between relational calculus and list comprehensions \cite{GHCListComprehension,ComprehensiveComprehensions} in order to avoid the significant theoretical performance hit.

It is widely noted that \emph{joins}, an integral operation in relational algebra are associated with inefficient implementations \cite{JoinProcessing}. It is easy to see why when considering the most general joins. However in their paper \cite{RelationalAlgebraByWayOfAdjunctions} they concern themselves with a specialised join called an \emph{equijoin}. As described in \fref{sec:joins} an equijoin is a specialised \emph{theta-join} -- a way of combining two relations based off of an arbitrary condition depending on the attributes of both relations. An integral part to the calculation of a theta-join is calculating the Cartesian product (all possible combinations of tuples of both relationships, as described in \fref{sec:products}). The algorithm must then filter every single tuple individually to check for equality of attributes! It is clear that this is wasteful for such an specialised join. As a more practical example consider the SQL program:
\begin{lstlisting}
  SELECT *
  FROM R, S
  WHERE R.a = S.b
\end{lstlisting}
This could be naively converted into a list comprehension with the following:
\[
  \left[\,(r, s)\;|\;r \leftarrow R,\;s \leftarrow S,\;r.a = s.b\,\right]
\]
where $(r, s)$ is seen as a single tuple whose attributes \attribute{\relation{R}.a} and \attribute{\relation{R}.b} are merged. \todo{fix example so that you do not need to see things}

 With this naive list comprehension implementation we effectively convert \equijoin{R}{\attribute{a}}{S}{\attribute{b}} to \select{\attribute{a} = \attribute{b}}{\relation{R} \times \relation{S}}, generating a relation with $|R||S|$ tuples in the process then filtering through each.

 This can much be much more efficiently implemented by viewing databases as indexed tables. We can index each relation by its associated attribute in the equijoin and merge the results -- localising the data required to in a cartesian product. This approach admits a linear time equijoin, if careful about comparison and projection functions. \cite{RelationalAlgebraByWayOfAdjunctions} With some mathematical tools explained in \fref{sec:gradedmonads} we can describe give these operations a monadic structure and therefore a comprehension syntax using the extended syntax discussed above.

 What this project adds to this story is a concrete demonstration of the improvement this solution offers.
 We will use the \emph{Haskell} and the list comprehensions and functions described above to implement a simple database querying software taking into account these key changes. Along with real world pragmatic data sources, benchmarking techniques found in \fref{sec:benchmarking} will be used to accurately measure and compare the efficiency difference between the two approaches -- with the new equijoin and without. This evidence would be very important to justify the use of these methods and the claims made in the paper \cite{RelationalAlgebraByWayOfAdjunctions}.
 With a concrete implementation, it could also provide insights into the downfall of remaining operations as well as those mentioned in the paper, potentially using profiling techniques to further analyse the performance bottlenecks, in order to inspire a theoretical approach at determining the issue as shown in \cite{RelationalAlgebraByWayOfAdjunctions} instead of an efficiency-driven optimisation approach.

\chapter{Background} % 10-20 pages
\begin{comment}
This should form the bulk of the interim report. You should consider that your objective here is to produce a near final version of the background section, as it will appear in your final report.  All of this material should be re-usable, so it is worth getting it right at this stage of the project.  The details of what to include can be found in the Project Report guidelines.

From the Project Report guidelines:
The background section of the report should set the project into context by relating it to existing published work which you read at the start of the project when your approach and methods were being considered. There are usually many ways of solving a given problem, and you shouldn't just pick one at random. Describe and evaluate as many alternative approaches as possible. The published work may be in the form of research papers found in the academic literature, articles, text books, technical manuals, or even existing software or hardware of which you have had hands-on experience. You must acknowledge the sources of your inspiration. You are expected to have seen and thought about other people's ideas; your contribution will be putting them into practice in some other context. However, avoid plagiarism: if you take another person's work as your own and do not cite your sources of information/inspiration you are being dishonest. When referring to other pieces of work, cite the sources where they are referred to or used, rather than just listing them at the end. Accidental plagiarism or not knowing how to cite and reference is not a valid reason for plagiarism. Make sure you read and digest the Department's plagiarism document .

In writing the Background chapter you must demonstrate your ability to *analyse*, *synthesise* and *apply critical* judgement. Analysis is shown by explaining how the proposed solution operates in your own words as well as its benefits and consequences. Synthesis is shown through the organisation of your Related Work section and through identifying and generalising common aspects across different solutions. Critical judgement is shown by discussing the limitations of the solutions proposed both in terms of their disadvantages and limits of applicability.

Typically you can look for Background work using different search engines including:
* Google Scholar
* IEEExplore
* ACM Digital Library
* Citeseer
* Science Direct

**Note 1:** Often the terms *Background, Literature Review, Related Work* and *State of the Art* are used interchangeably.
**Note 2**: Keyword search is wonderful, but you need the right *Keywords*.
**Note 3:** IEEExplore, ACM Digital Library and Science Direct may require you to be on the College network to download the PDF versions of papers. If at home, use VPN.
\end{comment}

\section{Crash course}
\subsection{Category theory}
\subsubsection{Categories}
\subsubsection{Functors}
\subsubsection{Natural transformations}
\subsubsection{Adjunctions}
\todo{Add more information between here}
\subsubsection{Graded Monads}

\subsection{Relational algebra}
\todo{Add citation to the database book}
\subsubsection{Projections}
\subsubsection{Selections}
\subsubsection{Products}
\subsubsection{Joins}

\subsection{Evolution of database representation}
\subsubsection{Bags}
\todo{Understand and distinguish between bags being the bulk type}
\paragraph{Characteristics of a database}We expect our database approximation to not be ordered and admit multiplicities and a finite bag of values is one of the simplest constructions that does so. Like a finite set, a bag contains a collection of unordered values. However, unlike a set, bags can contain duplicate elements. \cite{RelationalAlgebraByWayOfAdjunctions} This multiplicty is key for processing non-idempotent aggregations. For instance, if summing up the ages of a database of people, without admitting multiplicty we would only sum each unique age once.
\subparagraph{Generalisation}Furthermore, going forward we generalise to bags of any types instead of the classical ``bags of records''. This also allows us to deal with intermediate tables that contain non-record values.
\todo{Add the mathematical parts about bags}
\todo{When reading about fininte maps it says that it's better for databases as non finite maps cannot be aggreagated, can non finite bags be aggreagates - is this why they are finite?}

\subsubsection{Indexed tables}
We want to move towards an indexed representation of our table in order to equijoin by indexing. \todo{Understand if this is right and equijoin by indexing}. So in this section we introduce the mathematical concepts required to define such an implementation.
\theoremstyle{definition}\newtheorem*{psetdef}{Pointed set}
\theoremstyle{definition}\newtheorem*{ppfuncdef}{Point-preserving function}
\theoremstyle{definition}\newtheorem*{mapdef}{Map}
\theoremstyle{definition}\newtheorem*{finitemapdef}{Finite map}
\begin{psetdef}\label{def:pset}
  A pointed set $\pset{A}{a}$ is a set $A$ with a distinguished element $a \in A$.
\end{psetdef}
We commonly refer to the distinguished element of a set $A$ as $null_A$ or, when not ambiguous to do so, $null$.
\begin{ppfuncdef}\label{def:ppfunc}
  Given two pointed sets $\pset{A}{null_A}$ and $\pset{B}{null_B}$, a total function $f: A \rightarrow B$ is point-preserving if $f(null_A) = null_B$.
\end{ppfuncdef}
\todo{See if there is a way to have better spacing in brackets}

\todo{Add in more mathematical detail for point preserving functions}

We now have the mathematical tools required to define a map. In its finite form a map is widely known in computer science by many other names such as a dictionary, association lists or key-value maps.

Let $\keyset$ be a set and $\valset$ a pointed set. To those already familiar with maps, it may help to think of $\keyset$ as keys and $\valset$ as values.
\begin{mapdef}
  A map of type $\map{\keyset}{\valset}$ is a total function from K to V.
\end{mapdef}
\begin{finitemapdef}
  A finite map of type \finitemap{\keyset}{\valset} is a map where only a finite number of keys are mapped to $null_\valset$ (where $null_\valset$ is the distinguished element of \valset). 
\end{finitemapdef}
The advantage of using a finite map in a database is to allow aggregation.
\todo{Understand why only semimonoidal}
\todo{Introduced indexed tables}

\section{Literature Review}
\subsection{Benchmarking databases}
\subsection{Expanding the use of adjunctions}
\chapter{Project Plan} % 1-2 pages
\begin{comment}
You should explain what needs to be done in order to complete the project and roughly what you expect the timetable to be. Don’t forget to include the project write-up (the final report), as this is a major part of the exercise. It’s important to identify key milestones and also fall-back positions, in case you run out of time.  You should also identify what extensions could be added if time permits.  The plan should be complete and should include those parts that you have already addressed (make it clear how far you have progressed at the time of writing).  This material will *not* appear in the final report.
\end{comment}

The project has a clear starting point and a wide range of potential branches to eventually explore. This allows the project to venture down many different specific roots, ranging in application focused to theoretically based.
We outline the stages the project can take below.

\section{Initial research and literature review}\label{sec:researchstage}
As the project could eventually take a turn into the theoretical side of the paper \cite{RelationalAlgebraByWayOfAdjunctions} this part is crucial. The time frame for this will be ongoing and has started, however the bulk should be finished by the end of January. An intuition as well as knowledge of the applications of category theory in the wider range of computer science would be vital in order to spot potential improvements in any results found while implementing and benchmarking the findings of the base paper \cite{RelationalAlgebraByWayOfAdjunctions} .

I do expect, however, that there would be a period of time after finishing \fref{sec:implementation} to return to this step in more detail than normal in order to fully solidify the direction of the project.

Of course research is also important to create a fair and rigorous benchmarking system. This style of research should be concluded in the coming days in order to ensure that any implementation is built to play nicely with the tools to be used (and potentially initially specifically for the datasets chosen).
\section{Implementation of the differing equijoin algorithms}\label{sec:implementation}
After understanding the basics of the theoretical implications as well as how category theory underpins Haskell as whole, work can be started on the first tangible step. That is to implement a mini database system that can query a given dataset, both using the inefficient and efficient equijoins.
\subsection{Implementation}
Initially the implementation will be very tailored to the data sources chosen. This will be done to avoid the difficulties of generalising types and creating a coherent system under Haskell's strict rules. It is not too difficult or even out of the scope of the project, but I feel as though focusing on the results of benchmarking will be more efficient in earlier stages. I hope to get this done by the end of February.
\subsection{Evaluation and analysis}
The analysis should begin as soon as the implementation has been completed. At the very least, benchmarking seems to be quite heuristic and experience helps a lot, thus some time should be used to make sure that the results are as accurate as possible. At this stage, if experience has suggested that changing the implementation would not be too much effort, I might also consider changing the data source the tests were based on in order to deal with more interesting or accurate results if any issues seem to emerge. This process should hopefully take a couple weeks.

Evaluation, however, should be done as an ongoing process during the project. A plan to evaluate will be presented in \fref{sec:correctnessevaluation}.

\section{Theoretical analysis of the remaining relational algebra}\label{sec:theoreticalanalysis}
After the implementation of the findings of the paper, I hope to have gained a much deeper understanding of the practicalities of database querying. This along with the mathematical intuition I hope to have gained during the stage described \fref{sec:researchstage} would inspire thought to which other operations currently do not adhere to a monadic structure or efficient implementations. Of course this would be greatly aided by research in the large field of database queries, potentially choosing to describe already known optimisations using this new model and adjunctions instead. This section of the project could last until new work is no longer being found.
\section{Reporting the project}
% Report due Monday 19 June 2022
Reporting the project is a necessity. My project is unique in that the first section is itself an analysis and therefore, if done correctly, will almost immediately be written in a final report style while it is being conducted. Later parts of the project might need their own time set apart to write about as mathematics is rarely linear and the report should tell a compelling story.

My plan is to continue to develop the tools to efficiently write my report alongside the implementation of the database, and begin to draft and validate the final parts of the report a month before the deadline. Useful tools I have already integrated are both manual and commit triggered checks on the spelling and code style of the reports. I have also worked on automatic releases and management of different report versions, early in December. A few of the systems are helpful but not at their potential yet and so I might put some time to think about possible extensions, for instance custom dictionaries for the spell checker so it does not flag up domain specific terminology as errors. I definitely need to migrate the report structure I have to lhs2\TeX{} so that when I have code snippets to report it can be done with ease. This should be finished by February, in time to develop the database in the report if that is the workflow I find most beneficial.
\section{Possible extensions}
\subsection{Progress in optimising other aspects of relational algebra}
Given a successful stage in \fref{sec:theoreticalanalysis} this might give the project a lot of scope to grow as an extension over the final months. I could choose between using the category theory learned to write proofs of the theoretical advantage of any novel approach I think of, or just a novel description of what is already there. Alternatively I could extend my implementation to include the modified operations and do another analysis of their practical performance change.
\subsection{A full implementation}
If I choose to take this project down a more practical root, a full general implementation of the database system implemented could be made as an extension. This would allow a specific schema and interface with the data stored in Haskell's data structures itself, differing from database abstraction layers that communicate to SQL servers currently used in production.\cite{HDBC} Although this extension would not be production ready, I hope that it would be demonstrative on whether list comprehensions have a place among relational calculus, both by its viability to be implemented but also developer satisfaction using them.
\subsection{Applications of related fields to relational calculus}
Category theory is extremely general and has countless applications in computer science, from describing lambda calculi to describing type systems.\cite{BasicCategoryTheoryforCS} A particularly interesting paper about modelling concurrency \cite{CategoriesForModellingConcurrency} caught my eye as a topic to research into in the future, given how vital concurrency is to database management systems \cite{DatabaseSystems}.

Research could also be done into the viability of modelling mutating behaviours via referential calculus, that either change schema or components.

\chapter{Evaluation plan} % 1-2 pages
\begin{comment}
Project evaluation is very important, so it's important to think now about how you plan to measure success. For example, what functionality do you need to demonstrate?  What experiments to you need to undertake and what outcome(s) would constitute success?  What benchmarks should you use? How has your project extended the state of the art?  How do you measure qualitative aspects, such as ease of use?  These are the sort of questions that your project evaluation should address; this section should outline your plan.
\end{comment}
\section{During implementation}\label{sec:evaluationduringimplementation}
\chapter{Ethical issues} % 1-2 pages
\begin{comment}
What are the wider ethical, legal, professional and societal issues surrounding your project and the accompanying research? You should use the ethics checklist as the basis for this discussion. 
\end{comment}
The project itself (in its early stages) does not have a novel ideal as a useable output and thus is very neutral in its ethical impact. Instead it has the ability to recommend or withhold recommendation of an earlier approach to database optimisations via implementation and benchmarking results, and in its final stages may recommend further optimisations and so the more interesting ethical discussion is around the ethics of database optimisations.

Databases are a necessity for most modern applications. Furthermore, database optimisations are especially important for large, scaling products or services. This means that improving this area of computer science has a very indiscriminate effect on the type of software products that can make use of it. Of course, just as any private company could make use of the newer more efficient techniques, any military hosting a database could too. It is also worth noting the implementation is in a functional programming language, a rising paradigm \todo{cite} typically associated with reliability \todo{add another quality} \todo{Add citation for this} which aligns very well with the real-time and safety critical nature of military applications. \todo{cite that these are important for military applications}

In fact, it could be argued that any findings could benefit society at large in an environmental standpoint. The heart of the project is in database optimisations, which in turn reduce the amount of time required to process queries. \todo{research this} One could extrapolate, that without much of an increased processing complexity and resource consumption, this would lead to a reduction in the energy required for these database queries. Given how commonplace database queries are \todo{find stats} this could have a non-trivial impact.

It is worth noting that the early stages of the project involves implementing the ideas created in \cite{RelationalAlgebraByWayOfAdjunctions}. This causes no legal issues as there is no protection on an idea, lacking the tangibility (and ability to be patented) \todo{See exact laws}. On an ethical level, despite implementations, improvements and building on work is commonplace, encouraged and vital to the academic community, one of the main contributors of the paper is my supervisor and I interpret this as implicit consent. With good referencing practice, plagiarism would also be completely avoided.

There is an interesting discussion to be had on the ethical impact of the benchmarking dataset to be used in the early stages of the project. In order to evaluate the suggestions in \cite{RelationalAlgebraByWayOfAdjunctions} a database system will need to be implemented, but more consequentially one or more data sources will need to be used to evaluate the queries. I do not think that creating a novel data source will be very useful or worthwhile so I will need to acquire at least one somehow, raising at least two issues.
\subparagraph*{Data protection} The immediate issue that springs to mind would be whether the data source contains personal data and if so the subsequent data protection methods required to fulfil my legal and ethical obligation. To avoid this issue for primary data sources, I will not collect data when creating a benchmarking database, instead opting to randomly generate records. For data sources I acquire, they will all be in public domain already, and avoid sensitive personal data. This should minimise the risk to individuals affected. The data protection laws consider "organisation" and "structuring" as processing personal data, hence why this is relevant. The databases acquired might need to be reorganised into the data structured explained in this report \todo{cite}. \todo{See if this is enough} \todo{Cite data protection laws}
\subparagraph*{Database rights} Another consideration when outsourcing database collections, is the extent to which those who have created the database are protected. The creation of a database, although not creative, has been recognised as taking significant work and so those that create the database have a form of copyright protection on their work. \todo{Make sure that this is accurate, cite if so} This risk will be mitigated by ensuring that all databases used are in the public domain and the correct licenses are acquired.

\todo{Potentially comment on the distribution of my implementation}

\section{Extensions}
Furthermore, any domains investigated in this project have their own set of ethical considerations which reasonably cannot be discussed now (as they are largely unknown). Care should be taken on thinking of the ethical impact of the work that is being taken when branching out to other scenarios that use bulk-types. For instance, the suggested family of languages -- logic programming -- has close ties with artificial intelligence which itself is an ethically complex subject. \todo{Cite AI ethics} I think it is fair to say, however, that the work done will not bias any field in a particular way (ethical or not), but potentially contribute to more efficient implementations or different perspective of ideas within them.

% Back matter
\bibliographystyle{unsrtnat}
% DoC uses the Vancouver Referencing Format.
\bibliography{report/bibs/interim}

\end{document}
\newpage
\tableofcontents

% Main matter 
\documentclass[a4paper, twoside]{report}

%% Language and font encodings
\usepackage[english]{babel}
\usepackage[utf8x]{inputenc}
\usepackage[T1]{fontenc}
\usepackage{mathrsfs}

%% Sets page size and margins
\usepackage[a4paper,top=3cm,bottom=2cm,left=3cm,right=3cm,marginparwidth=1.75cm]{geometry}

\usepackage{amsthm}
\usepackage{graphicx}
\usepackage[]{listings}
\usepackage{booktabs}
\usepackage[colorinlistoftodos]{todonotes}
\usepackage[colorlinks=true, allcolors=blue]{hyperref}
\usepackage{comment}
\usepackage[numbers]{natbib}
\usepackage{fancyref}
\usepackage{wrapfig}

%% Math
\usepackage{mathtools}
\usepackage{stmaryrd}

\usepackage{title/titlestructure}

%% Utils
\usepackage{background/utils}


\title{A Deeper Dive into Relational Algebra by Way of Adjunctions}
\author{Matteo Bongiovanni}
\date{January 2023}
\supervisor{Dr.\ Nicolas Wu}
\secondmarker{Dr.\ Steffen van Bakel}
\documenttype{MEng Individual Project}
\university{Imperial College London}
\department{Department of Computing}
\degree{Joint Mathematics and Computing}

\begin{document}
\begin{comment}
General Information
The project report (occasionally referred to as a thesis) is an extremely important aspect of the project. It serves to show what you have achieved and should demonstrate that:

You understand the wider context of computing by relating your choice of project, and the approach you take, to existing products or research.
You can apply the theoretical and practical techniques taught in the course to the problem you are addressing and that you understand their relevance to the wider world of computing.
You are capable of objectively criticising your own work and making constructive suggestions for improvements or further work based on your experiences so far.
As a computing professional, you operate ethically and can explain your thinking and working processes clearly and concisely to third parties who may not be experts in the field in which you are working.
Most of the project assessors will not have followed the project throughout and will only have a short time to listen to a presentation or see a demonstration. For this reason they will rely heavily on the report to judge the project.

Many students underestimate the importance of the report and make the mistake of thinking that top marks can be achieved simply for producing a good product. This is fundamentally not the case and many projects have been graded well below their potential because of an indifferent or poor write-up. In order to get the balance right you should consider that the aim of the project is to produce a good report and that software, hardware, theory etc. that you developed during the project are merely a means to this end. Don't make the mistake of leaving the write-up to the last minute. Ideally you should produce the bulk of the report as you go along and use the last week or two to bring it together into a coherent document.

It is helpful to get feedback from your supervisor about your project report, but supervisors cannot be expected to look at documents written at the last minute or at more than one chapter at a time. Allow plenty of time for this.

Project Report Length
The nominal maximum report length is 60 pages, plus appendices. Shorter reports may be appropriate for many projects and longer reports can be submitted with prior approval (see below).

The page count includes:

title page
abstract
dedications
table of contents
The page count does not include:

bibliography
appendices (but the report should make sense without the appendices)
When writing up your work be reminded of the fact that a 60-page report isn't necessarily better than a 30-page report, in the same way that a 10,000 line implementation isn't necessarily better 5,000 line one. Conciseness, clarity and elegance are invaluable qualities in report writing, just as they are in programming. Indeed, you will lose credit if it is evident to the assessors that the report has been padded with superfluous content in order to make it appear more substantial than the work justifies. The assessors greatly value quality over quantity.

It is important to understand that the report is assessed on the quality of the technical writing. Here this means your ability to document your ideas, design, implementation, evaluation, conclusions etc. succinctly, using appropriate scientific language, as might a researcher when writing up their findings in an technical magazine, academic journal or conference proceedings.

If you have good reason to use more than 60 pages you should consult your supervisor and second marker and obtain their approval in advance - they will be required to read and assess the whole report, whereas other assessment team members may only scan key sections.

Project Layout
The physical layout and formatting of the report is also important. A tidy, well laid out and consistently formatted document makes for easier reading and is suggestive of a careful and professional attitude towards its preparation. Both LaTeX and MS Word will allow you to produce a cleanly formatted document. Many supervisors will advise you to use LaTeX as this solves most of the formatting problems for you. If your report is to contain a substantial number of mathematical formulae you are strongly advised to use LaTeX.

Project Report Structure
The exact structure of your project (the chapter titles etc) is up to you. Discuss with your supervisor what structure works best for your project.

A project report will usually contain a number of chapters (Chapter 1, Chapter 2 etc). Each chapter contains Sections (2.1, 2.2 etc) and may also contain subsections (2.1.1, 2.1.2) etc. Try to avoid too many levels of subheading - three is usually sufficient.

Projects will usually contain all of the following elements, although these may not exactly map to chapters.
\end{comment}

% Front matter
\begin{comment}
This should include the project title and the name of the author of the report. You can also list the name of your supervisor if you wish. 
\end{comment}

\begin{titlepage}

\newcommand{\HRule}{\rule{\linewidth}{0.5mm}} % Defines a new command for the horizontal lines, change thickness here

%----------------------------------------------------------------------------------------
%	LOGO SECTION
%----------------------------------------------------------------------------------------

\includegraphics[width=8cm, draft=false]{report/title/logo.eps}\\[1cm] % Include a department/university logo - this will require the graphicx package
 
%----------------------------------------------------------------------------------------

\center\ % Center everything on the page

%----------------------------------------------------------------------------------------
%	HEADING SECTIONS
%----------------------------------------------------------------------------------------

\makeatletter
\textsc{\LARGE \@documenttype}\\[1.5cm]
\textsc{\Large \@university}\\[0.5cm]
\ifdefined\@degree\
\textsc{\large \@degree}\\[0.5cm]
\fi
\textsc{\large \@department}\\[0.5cm]
\makeatother

%----------------------------------------------------------------------------------------
%	TITLE SECTION
%----------------------------------------------------------------------------------------
\makeatletter
\HRule\\[0.4cm]
{ \huge \bfseries \@title}\\[0.4cm] % Title of your document
\HRule\\[1.5cm]
 
%----------------------------------------------------------------------------------------
%	AUTHOR SECTION
%----------------------------------------------------------------------------------------

\begin{minipage}{0.4\textwidth}
\begin{flushleft} \large
\emph{Author:}\\
\@author\ % Your name
\end{flushleft}
\end{minipage}
~\begin{minipage}{0.4\textwidth}
\begin{flushright} \large
\emph{Supervisor:} \\
\@supervisor\\[1.2em]

\ifdefined\@secondmarker\
\emph{Second Marker:} \\
\@secondmarker\
\fi

\end{flushright}
\end{minipage}\\[2cm]

%----------------------------------------------------------------------------------------
%	DATE SECTION
%----------------------------------------------------------------------------------------

{\large \@date}\\[2cm] % Date, change the \today to a set date if you want to be precise

\vfill % Fill the rest of the page with whitespace
\makeatother

\end{titlepage}
\begin{comment}
The abstract is a very brief summary of the report's contents. It should be no more than half a page long. Somebody unfamiliar with your project should have a good idea of what it's about having read the abstract alone and will know whether it will be of interest to them. Note that the abstract is a summary of the entire project including its conclusions. A common mistake is to provide only introductory elements in the abstract without saying what has been achieved.
\end{comment}


\begin{abstract}
In order to evaluate alternative equijoin implementation in Haskell, this paper
introduced the \relation{JOINBENCH} relation and surrounding methodology and
tooling. The \relation{JOINBENCH} relation is a scheme designed for the
synthesis of data optimised to help evaluate the performance of equijoins in a
variety of scenarios. Furthermore, this paper presents a low-level library that
helps users customise and define their own synthetic data sources for future
benchmarking purposes.

In their distinguished paper
\relalg{}~\cite{RelationalAlgebraByWayOfAdjunctions} it was noted that the
monadic structure of bulk types can help explain most of relational algebra.
Using this structure, the authors designed a new method to facilitate
the use of monad comprehensions in an efficient implementation of equijoins of
relational databases. This project presents an implementation of such a system
and an evaluation of the performance such query optimisations carry using the
bespoke tooling described above. 
\end{abstract}

\tableofcontents

% Main matter 
\chapter{Introduction} % 1-3 pages
\begin{comment}
It’s a good idea to *try* to write the introduction to your final report early on in the project. However, you will find it hard, as you won’t yet have a complete story and you won’t know what your main contributions are going to be. However, the exercise is useful as it will tell you what you *don’t* yet know and thus what questions your project should aim to answer. For the interim report this section should be a short, succinct, summary of the project’s main objectives. Some of this material may be re-usable in your final report, but the chances are that your final introduction will be quite different.  You are therefore advised to keep this part of the interim report short, focusing on the following questions: What is the problem, why is it interesting and what’s your main idea for solving it?  (DON'T use those three questions as subheadings however!  The answers should emerge from what you write.)
\end{comment}

Databases are absolutely vital to modern day society and contain domain specific knowledge about anything from specialised images of eyes \todo{add citation here} to the structure of crystals \cite{CambridgeStructuralDatabase}. Since its conception many different data models describing how to hold the data in databases have emerged, including the relational model \cite{RelationalModel} and the semi-structured model\cite{DatabaseSystems}.

In this project we concern ourselves with the relational model (to be introduced in \fref{sec:relationalmodel}) as a way of modelling the database. This rich model has many methods of expressing queries, especially relational algebra and relational calculus, both with their strengths and weaknesses \cite{RelationalCalculus,RelationalModel}. \todo{they are equivalent though?} However, the favoured specification of the authors of \cite{RelationalAlgebraByWayOfAdjunctions} seemed to be through list comprehensions, a beautiful feature that ``provide for a concise and expressive notation for writing list-processing code''. \cite{MonadComprehensions} They eventually propose using GHC's extended list comprehension syntax specifically designed to help bridge the already close relationship between relational calculus and list comprehensions \cite{GHCListComprehension,ComprehensiveComprehensions} in order to avoid the significant theoretical performance hit.

It is widely noted that \emph{joins}, an integral operation in relational algebra are associated with inefficient implementations \cite{JoinProcessing}. It is easy to see why when considering the most general joins. However in their paper \cite{RelationalAlgebraByWayOfAdjunctions} they concern themselves with a specialised join called an \emph{equijoin}. As described in \fref{sec:joins} an equijoin is a specialised \emph{theta-join} -- a way of combining two relations based off of an arbitrary condition depending on the attributes of both relations. An integral part to the calculation of a theta-join is calculating the Cartesian product (all possible combinations of tuples of both relationships, as described in \fref{sec:products}). The algorithm must then filter every single tuple individually to check for equality of attributes! It is clear that this is wasteful for such an specialised join. As a more practical example consider the SQL program:
\begin{lstlisting}
  SELECT *
  FROM R, S
  WHERE R.a = S.b
\end{lstlisting}
This could be naively converted into a list comprehension with the following:
\[
  \left[\,(r, s)\;|\;r \leftarrow R,\;s \leftarrow S,\;r.a = s.b\,\right]
\]
where $(r, s)$ is seen as a single tuple whose attributes \attribute{\relation{R}.a} and \attribute{\relation{R}.b} are merged. \todo{fix example so that you do not need to see things}

 With this naive list comprehension implementation we effectively convert \equijoin{R}{\attribute{a}}{S}{\attribute{b}} to \select{\attribute{a} = \attribute{b}}{\relation{R} \times \relation{S}}, generating a relation with $|R||S|$ tuples in the process then filtering through each.

 This can much be much more efficiently implemented by viewing databases as indexed tables. We can index each relation by its associated attribute in the equijoin and merge the results -- localising the data required to in a cartesian product. This approach admits a linear time equijoin, if careful about comparison and projection functions. \cite{RelationalAlgebraByWayOfAdjunctions} With some mathematical tools explained in \fref{sec:gradedmonads} we can describe give these operations a monadic structure and therefore a comprehension syntax using the extended syntax discussed above.

 What this project adds to this story is a concrete demonstration of the improvement this solution offers.
 We will use the \emph{Haskell} and the list comprehensions and functions described above to implement a simple database querying software taking into account these key changes. Along with real world pragmatic data sources, benchmarking techniques found in \fref{sec:benchmarking} will be used to accurately measure and compare the efficiency difference between the two approaches -- with the new equijoin and without. This evidence would be very important to justify the use of these methods and the claims made in the paper \cite{RelationalAlgebraByWayOfAdjunctions}.
 With a concrete implementation, it could also provide insights into the downfall of remaining operations as well as those mentioned in the paper, potentially using profiling techniques to further analyse the performance bottlenecks, in order to inspire a theoretical approach at determining the issue as shown in \cite{RelationalAlgebraByWayOfAdjunctions} instead of an efficiency-driven optimisation approach.

\chapter{Background} % 10-20 pages
\begin{comment}
This should form the bulk of the interim report. You should consider that your objective here is to produce a near final version of the background section, as it will appear in your final report.  All of this material should be re-usable, so it is worth getting it right at this stage of the project.  The details of what to include can be found in the Project Report guidelines.

From the Project Report guidelines:
The background section of the report should set the project into context by relating it to existing published work which you read at the start of the project when your approach and methods were being considered. There are usually many ways of solving a given problem, and you shouldn't just pick one at random. Describe and evaluate as many alternative approaches as possible. The published work may be in the form of research papers found in the academic literature, articles, text books, technical manuals, or even existing software or hardware of which you have had hands-on experience. You must acknowledge the sources of your inspiration. You are expected to have seen and thought about other people's ideas; your contribution will be putting them into practice in some other context. However, avoid plagiarism: if you take another person's work as your own and do not cite your sources of information/inspiration you are being dishonest. When referring to other pieces of work, cite the sources where they are referred to or used, rather than just listing them at the end. Accidental plagiarism or not knowing how to cite and reference is not a valid reason for plagiarism. Make sure you read and digest the Department's plagiarism document .

In writing the Background chapter you must demonstrate your ability to *analyse*, *synthesise* and *apply critical* judgement. Analysis is shown by explaining how the proposed solution operates in your own words as well as its benefits and consequences. Synthesis is shown through the organisation of your Related Work section and through identifying and generalising common aspects across different solutions. Critical judgement is shown by discussing the limitations of the solutions proposed both in terms of their disadvantages and limits of applicability.

Typically you can look for Background work using different search engines including:
* Google Scholar
* IEEExplore
* ACM Digital Library
* Citeseer
* Science Direct

**Note 1:** Often the terms *Background, Literature Review, Related Work* and *State of the Art* are used interchangeably.
**Note 2**: Keyword search is wonderful, but you need the right *Keywords*.
**Note 3:** IEEExplore, ACM Digital Library and Science Direct may require you to be on the College network to download the PDF versions of papers. If at home, use VPN.
\end{comment}

\section{Crash course}
\subsection{Category theory}
\subsubsection{Categories}
\subsubsection{Functors}
\subsubsection{Natural transformations}
\subsubsection{Adjunctions}
\todo{Add more information between here}
\subsubsection{Graded Monads}

\subsection{Relational algebra}
\todo{Add citation to the database book}
\subsubsection{Projections}
\subsubsection{Selections}
\subsubsection{Products}
\subsubsection{Joins}

\subsection{Evolution of database representation}
\subsubsection{Bags}
\todo{Understand and distinguish between bags being the bulk type}
\paragraph{Characteristics of a database}We expect our database approximation to not be ordered and admit multiplicities and a finite bag of values is one of the simplest constructions that does so. Like a finite set, a bag contains a collection of unordered values. However, unlike a set, bags can contain duplicate elements. \cite{RelationalAlgebraByWayOfAdjunctions} This multiplicty is key for processing non-idempotent aggregations. For instance, if summing up the ages of a database of people, without admitting multiplicty we would only sum each unique age once.
\subparagraph{Generalisation}Furthermore, going forward we generalise to bags of any types instead of the classical ``bags of records''. This also allows us to deal with intermediate tables that contain non-record values.
\todo{Add the mathematical parts about bags}
\todo{When reading about fininte maps it says that it's better for databases as non finite maps cannot be aggreagated, can non finite bags be aggreagates - is this why they are finite?}

\subsubsection{Indexed tables}
We want to move towards an indexed representation of our table in order to equijoin by indexing. \todo{Understand if this is right and equijoin by indexing}. So in this section we introduce the mathematical concepts required to define such an implementation.
\theoremstyle{definition}\newtheorem*{psetdef}{Pointed set}
\theoremstyle{definition}\newtheorem*{ppfuncdef}{Point-preserving function}
\theoremstyle{definition}\newtheorem*{mapdef}{Map}
\theoremstyle{definition}\newtheorem*{finitemapdef}{Finite map}
\begin{psetdef}\label{def:pset}
  A pointed set $\pset{A}{a}$ is a set $A$ with a distinguished element $a \in A$.
\end{psetdef}
We commonly refer to the distinguished element of a set $A$ as $null_A$ or, when not ambiguous to do so, $null$.
\begin{ppfuncdef}\label{def:ppfunc}
  Given two pointed sets $\pset{A}{null_A}$ and $\pset{B}{null_B}$, a total function $f: A \rightarrow B$ is point-preserving if $f(null_A) = null_B$.
\end{ppfuncdef}
\todo{See if there is a way to have better spacing in brackets}

\todo{Add in more mathematical detail for point preserving functions}

We now have the mathematical tools required to define a map. In its finite form a map is widely known in computer science by many other names such as a dictionary, association lists or key-value maps.

Let $\keyset$ be a set and $\valset$ a pointed set. To those already familiar with maps, it may help to think of $\keyset$ as keys and $\valset$ as values.
\begin{mapdef}
  A map of type $\map{\keyset}{\valset}$ is a total function from K to V.
\end{mapdef}
\begin{finitemapdef}
  A finite map of type \finitemap{\keyset}{\valset} is a map where only a finite number of keys are mapped to $null_\valset$ (where $null_\valset$ is the distinguished element of \valset). 
\end{finitemapdef}
The advantage of using a finite map in a database is to allow aggregation.
\todo{Understand why only semimonoidal}
\todo{Introduced indexed tables}

\section{Literature Review}
\subsection{Benchmarking databases}
\subsection{Expanding the use of adjunctions}
\chapter{Project Plan} % 1-2 pages
\begin{comment}
You should explain what needs to be done in order to complete the project and roughly what you expect the timetable to be. Don’t forget to include the project write-up (the final report), as this is a major part of the exercise. It’s important to identify key milestones and also fall-back positions, in case you run out of time.  You should also identify what extensions could be added if time permits.  The plan should be complete and should include those parts that you have already addressed (make it clear how far you have progressed at the time of writing).  This material will *not* appear in the final report.
\end{comment}

The project has a clear starting point and a wide range of potential branches to eventually explore. This allows the project to venture down many different specific roots, ranging in application focused to theoretically based.
We outline the stages the project can take below.

\section{Initial research and literature review}\label{sec:researchstage}
As the project could eventually take a turn into the theoretical side of the paper \cite{RelationalAlgebraByWayOfAdjunctions} this part is crucial. The time frame for this will be ongoing and has started, however the bulk should be finished by the end of January. An intuition as well as knowledge of the applications of category theory in the wider range of computer science would be vital in order to spot potential improvements in any results found while implementing and benchmarking the findings of the base paper \cite{RelationalAlgebraByWayOfAdjunctions} .

I do expect, however, that there would be a period of time after finishing \fref{sec:implementation} to return to this step in more detail than normal in order to fully solidify the direction of the project.

Of course research is also important to create a fair and rigorous benchmarking system. This style of research should be concluded in the coming days in order to ensure that any implementation is built to play nicely with the tools to be used (and potentially initially specifically for the datasets chosen).
\section{Implementation of the differing equijoin algorithms}\label{sec:implementation}
After understanding the basics of the theoretical implications as well as how category theory underpins Haskell as whole, work can be started on the first tangible step. That is to implement a mini database system that can query a given dataset, both using the inefficient and efficient equijoins.
\subsection{Implementation}
Initially the implementation will be very tailored to the data sources chosen. This will be done to avoid the difficulties of generalising types and creating a coherent system under Haskell's strict rules. It is not too difficult or even out of the scope of the project, but I feel as though focusing on the results of benchmarking will be more efficient in earlier stages. I hope to get this done by the end of February.
\subsection{Evaluation and analysis}
The analysis should begin as soon as the implementation has been completed. At the very least, benchmarking seems to be quite heuristic and experience helps a lot, thus some time should be used to make sure that the results are as accurate as possible. At this stage, if experience has suggested that changing the implementation would not be too much effort, I might also consider changing the data source the tests were based on in order to deal with more interesting or accurate results if any issues seem to emerge. This process should hopefully take a couple weeks.

Evaluation, however, should be done as an ongoing process during the project. A plan to evaluate will be presented in \fref{sec:correctnessevaluation}.

\section{Theoretical analysis of the remaining relational algebra}\label{sec:theoreticalanalysis}
After the implementation of the findings of the paper, I hope to have gained a much deeper understanding of the practicalities of database querying. This along with the mathematical intuition I hope to have gained during the stage described \fref{sec:researchstage} would inspire thought to which other operations currently do not adhere to a monadic structure or efficient implementations. Of course this would be greatly aided by research in the large field of database queries, potentially choosing to describe already known optimisations using this new model and adjunctions instead. This section of the project could last until new work is no longer being found.
\section{Reporting the project}
% Report due Monday 19 June 2022
Reporting the project is a necessity. My project is unique in that the first section is itself an analysis and therefore, if done correctly, will almost immediately be written in a final report style while it is being conducted. Later parts of the project might need their own time set apart to write about as mathematics is rarely linear and the report should tell a compelling story.

My plan is to continue to develop the tools to efficiently write my report alongside the implementation of the database, and begin to draft and validate the final parts of the report a month before the deadline. Useful tools I have already integrated are both manual and commit triggered checks on the spelling and code style of the reports. I have also worked on automatic releases and management of different report versions, early in December. A few of the systems are helpful but not at their potential yet and so I might put some time to think about possible extensions, for instance custom dictionaries for the spell checker so it does not flag up domain specific terminology as errors. I definitely need to migrate the report structure I have to lhs2\TeX{} so that when I have code snippets to report it can be done with ease. This should be finished by February, in time to develop the database in the report if that is the workflow I find most beneficial.
\section{Possible extensions}
\subsection{Progress in optimising other aspects of relational algebra}
Given a successful stage in \fref{sec:theoreticalanalysis} this might give the project a lot of scope to grow as an extension over the final months. I could choose between using the category theory learned to write proofs of the theoretical advantage of any novel approach I think of, or just a novel description of what is already there. Alternatively I could extend my implementation to include the modified operations and do another analysis of their practical performance change.
\subsection{A full implementation}
If I choose to take this project down a more practical root, a full general implementation of the database system implemented could be made as an extension. This would allow a specific schema and interface with the data stored in Haskell's data structures itself, differing from database abstraction layers that communicate to SQL servers currently used in production.\cite{HDBC} Although this extension would not be production ready, I hope that it would be demonstrative on whether list comprehensions have a place among relational calculus, both by its viability to be implemented but also developer satisfaction using them.
\subsection{Applications of related fields to relational calculus}
Category theory is extremely general and has countless applications in computer science, from describing lambda calculi to describing type systems.\cite{BasicCategoryTheoryforCS} A particularly interesting paper about modelling concurrency \cite{CategoriesForModellingConcurrency} caught my eye as a topic to research into in the future, given how vital concurrency is to database management systems \cite{DatabaseSystems}.

Research could also be done into the viability of modelling mutating behaviours via referential calculus, that either change schema or components.

\chapter{Evaluation plan} % 1-2 pages
\begin{comment}
Project evaluation is very important, so it's important to think now about how you plan to measure success. For example, what functionality do you need to demonstrate?  What experiments to you need to undertake and what outcome(s) would constitute success?  What benchmarks should you use? How has your project extended the state of the art?  How do you measure qualitative aspects, such as ease of use?  These are the sort of questions that your project evaluation should address; this section should outline your plan.
\end{comment}
\section{During implementation}\label{sec:evaluationduringimplementation}
\chapter{Ethical issues} % 1-2 pages
\begin{comment}
What are the wider ethical, legal, professional and societal issues surrounding your project and the accompanying research? You should use the ethics checklist as the basis for this discussion. 
\end{comment}
The project itself (in its early stages) does not have a novel ideal as a useable output and thus is very neutral in its ethical impact. Instead it has the ability to recommend or withhold recommendation of an earlier approach to database optimisations via implementation and benchmarking results, and in its final stages may recommend further optimisations and so the more interesting ethical discussion is around the ethics of database optimisations.

Databases are a necessity for most modern applications. Furthermore, database optimisations are especially important for large, scaling products or services. This means that improving this area of computer science has a very indiscriminate effect on the type of software products that can make use of it. Of course, just as any private company could make use of the newer more efficient techniques, any military hosting a database could too. It is also worth noting the implementation is in a functional programming language, a rising paradigm \todo{cite} typically associated with reliability \todo{add another quality} \todo{Add citation for this} which aligns very well with the real-time and safety critical nature of military applications. \todo{cite that these are important for military applications}

In fact, it could be argued that any findings could benefit society at large in an environmental standpoint. The heart of the project is in database optimisations, which in turn reduce the amount of time required to process queries. \todo{research this} One could extrapolate, that without much of an increased processing complexity and resource consumption, this would lead to a reduction in the energy required for these database queries. Given how commonplace database queries are \todo{find stats} this could have a non-trivial impact.

It is worth noting that the early stages of the project involves implementing the ideas created in \cite{RelationalAlgebraByWayOfAdjunctions}. This causes no legal issues as there is no protection on an idea, lacking the tangibility (and ability to be patented) \todo{See exact laws}. On an ethical level, despite implementations, improvements and building on work is commonplace, encouraged and vital to the academic community, one of the main contributors of the paper is my supervisor and I interpret this as implicit consent. With good referencing practice, plagiarism would also be completely avoided.

There is an interesting discussion to be had on the ethical impact of the benchmarking dataset to be used in the early stages of the project. In order to evaluate the suggestions in \cite{RelationalAlgebraByWayOfAdjunctions} a database system will need to be implemented, but more consequentially one or more data sources will need to be used to evaluate the queries. I do not think that creating a novel data source will be very useful or worthwhile so I will need to acquire at least one somehow, raising at least two issues.
\subparagraph*{Data protection} The immediate issue that springs to mind would be whether the data source contains personal data and if so the subsequent data protection methods required to fulfil my legal and ethical obligation. To avoid this issue for primary data sources, I will not collect data when creating a benchmarking database, instead opting to randomly generate records. For data sources I acquire, they will all be in public domain already, and avoid sensitive personal data. This should minimise the risk to individuals affected. The data protection laws consider "organisation" and "structuring" as processing personal data, hence why this is relevant. The databases acquired might need to be reorganised into the data structured explained in this report \todo{cite}. \todo{See if this is enough} \todo{Cite data protection laws}
\subparagraph*{Database rights} Another consideration when outsourcing database collections, is the extent to which those who have created the database are protected. The creation of a database, although not creative, has been recognised as taking significant work and so those that create the database have a form of copyright protection on their work. \todo{Make sure that this is accurate, cite if so} This risk will be mitigated by ensuring that all databases used are in the public domain and the correct licenses are acquired.

\todo{Potentially comment on the distribution of my implementation}

\section{Extensions}
Furthermore, any domains investigated in this project have their own set of ethical considerations which reasonably cannot be discussed now (as they are largely unknown). Care should be taken on thinking of the ethical impact of the work that is being taken when branching out to other scenarios that use bulk-types. For instance, the suggested family of languages -- logic programming -- has close ties with artificial intelligence which itself is an ethically complex subject. \todo{Cite AI ethics} I think it is fair to say, however, that the work done will not bias any field in a particular way (ethical or not), but potentially contribute to more efficient implementations or different perspective of ideas within them.

% Back matter
\bibliographystyle{unsrtnat}
% DoC uses the Vancouver Referencing Format.
\bibliography{report/bibs/interim}

\end{document}
\documentclass[a4paper, twoside]{report}

%% Language and font encodings
\usepackage[english]{babel}
\usepackage[utf8x]{inputenc}
\usepackage[T1]{fontenc}
\usepackage{mathrsfs}

%% Sets page size and margins
\usepackage[a4paper,top=3cm,bottom=2cm,left=3cm,right=3cm,marginparwidth=1.75cm]{geometry}

\usepackage{amsthm}
\usepackage{graphicx}
\usepackage[]{listings}
\usepackage{booktabs}
\usepackage[colorinlistoftodos]{todonotes}
\usepackage[colorlinks=true, allcolors=blue]{hyperref}
\usepackage{comment}
\usepackage[numbers]{natbib}
\usepackage{fancyref}
\usepackage{wrapfig}

%% Math
\usepackage{mathtools}
\usepackage{stmaryrd}

\usepackage{title/titlestructure}

%% Utils
\usepackage{background/utils}


\title{A Deeper Dive into Relational Algebra by Way of Adjunctions}
\author{Matteo Bongiovanni}
\date{January 2023}
\supervisor{Dr.\ Nicolas Wu}
\secondmarker{Dr.\ Steffen van Bakel}
\documenttype{MEng Individual Project}
\university{Imperial College London}
\department{Department of Computing}
\degree{Joint Mathematics and Computing}

\begin{document}
\begin{comment}
General Information
The project report (occasionally referred to as a thesis) is an extremely important aspect of the project. It serves to show what you have achieved and should demonstrate that:

You understand the wider context of computing by relating your choice of project, and the approach you take, to existing products or research.
You can apply the theoretical and practical techniques taught in the course to the problem you are addressing and that you understand their relevance to the wider world of computing.
You are capable of objectively criticising your own work and making constructive suggestions for improvements or further work based on your experiences so far.
As a computing professional, you operate ethically and can explain your thinking and working processes clearly and concisely to third parties who may not be experts in the field in which you are working.
Most of the project assessors will not have followed the project throughout and will only have a short time to listen to a presentation or see a demonstration. For this reason they will rely heavily on the report to judge the project.

Many students underestimate the importance of the report and make the mistake of thinking that top marks can be achieved simply for producing a good product. This is fundamentally not the case and many projects have been graded well below their potential because of an indifferent or poor write-up. In order to get the balance right you should consider that the aim of the project is to produce a good report and that software, hardware, theory etc. that you developed during the project are merely a means to this end. Don't make the mistake of leaving the write-up to the last minute. Ideally you should produce the bulk of the report as you go along and use the last week or two to bring it together into a coherent document.

It is helpful to get feedback from your supervisor about your project report, but supervisors cannot be expected to look at documents written at the last minute or at more than one chapter at a time. Allow plenty of time for this.

Project Report Length
The nominal maximum report length is 60 pages, plus appendices. Shorter reports may be appropriate for many projects and longer reports can be submitted with prior approval (see below).

The page count includes:

title page
abstract
dedications
table of contents
The page count does not include:

bibliography
appendices (but the report should make sense without the appendices)
When writing up your work be reminded of the fact that a 60-page report isn't necessarily better than a 30-page report, in the same way that a 10,000 line implementation isn't necessarily better 5,000 line one. Conciseness, clarity and elegance are invaluable qualities in report writing, just as they are in programming. Indeed, you will lose credit if it is evident to the assessors that the report has been padded with superfluous content in order to make it appear more substantial than the work justifies. The assessors greatly value quality over quantity.

It is important to understand that the report is assessed on the quality of the technical writing. Here this means your ability to document your ideas, design, implementation, evaluation, conclusions etc. succinctly, using appropriate scientific language, as might a researcher when writing up their findings in an technical magazine, academic journal or conference proceedings.

If you have good reason to use more than 60 pages you should consult your supervisor and second marker and obtain their approval in advance - they will be required to read and assess the whole report, whereas other assessment team members may only scan key sections.

Project Layout
The physical layout and formatting of the report is also important. A tidy, well laid out and consistently formatted document makes for easier reading and is suggestive of a careful and professional attitude towards its preparation. Both LaTeX and MS Word will allow you to produce a cleanly formatted document. Many supervisors will advise you to use LaTeX as this solves most of the formatting problems for you. If your report is to contain a substantial number of mathematical formulae you are strongly advised to use LaTeX.

Project Report Structure
The exact structure of your project (the chapter titles etc) is up to you. Discuss with your supervisor what structure works best for your project.

A project report will usually contain a number of chapters (Chapter 1, Chapter 2 etc). Each chapter contains Sections (2.1, 2.2 etc) and may also contain subsections (2.1.1, 2.1.2) etc. Try to avoid too many levels of subheading - three is usually sufficient.

Projects will usually contain all of the following elements, although these may not exactly map to chapters.
\end{comment}

% Front matter
\begin{comment}
This should include the project title and the name of the author of the report. You can also list the name of your supervisor if you wish. 
\end{comment}

\begin{titlepage}

\newcommand{\HRule}{\rule{\linewidth}{0.5mm}} % Defines a new command for the horizontal lines, change thickness here

%----------------------------------------------------------------------------------------
%	LOGO SECTION
%----------------------------------------------------------------------------------------

\includegraphics[width=8cm, draft=false]{report/title/logo.eps}\\[1cm] % Include a department/university logo - this will require the graphicx package
 
%----------------------------------------------------------------------------------------

\center\ % Center everything on the page

%----------------------------------------------------------------------------------------
%	HEADING SECTIONS
%----------------------------------------------------------------------------------------

\makeatletter
\textsc{\LARGE \@documenttype}\\[1.5cm]
\textsc{\Large \@university}\\[0.5cm]
\ifdefined\@degree\
\textsc{\large \@degree}\\[0.5cm]
\fi
\textsc{\large \@department}\\[0.5cm]
\makeatother

%----------------------------------------------------------------------------------------
%	TITLE SECTION
%----------------------------------------------------------------------------------------
\makeatletter
\HRule\\[0.4cm]
{ \huge \bfseries \@title}\\[0.4cm] % Title of your document
\HRule\\[1.5cm]
 
%----------------------------------------------------------------------------------------
%	AUTHOR SECTION
%----------------------------------------------------------------------------------------

\begin{minipage}{0.4\textwidth}
\begin{flushleft} \large
\emph{Author:}\\
\@author\ % Your name
\end{flushleft}
\end{minipage}
~\begin{minipage}{0.4\textwidth}
\begin{flushright} \large
\emph{Supervisor:} \\
\@supervisor\\[1.2em]

\ifdefined\@secondmarker\
\emph{Second Marker:} \\
\@secondmarker\
\fi

\end{flushright}
\end{minipage}\\[2cm]

%----------------------------------------------------------------------------------------
%	DATE SECTION
%----------------------------------------------------------------------------------------

{\large \@date}\\[2cm] % Date, change the \today to a set date if you want to be precise

\vfill % Fill the rest of the page with whitespace
\makeatother

\end{titlepage}
\begin{comment}
The abstract is a very brief summary of the report's contents. It should be no more than half a page long. Somebody unfamiliar with your project should have a good idea of what it's about having read the abstract alone and will know whether it will be of interest to them. Note that the abstract is a summary of the entire project including its conclusions. A common mistake is to provide only introductory elements in the abstract without saying what has been achieved.
\end{comment}


\begin{abstract}
In order to evaluate alternative equijoin implementation in Haskell, this paper
introduced the \relation{JOINBENCH} relation and surrounding methodology and
tooling. The \relation{JOINBENCH} relation is a scheme designed for the
synthesis of data optimised to help evaluate the performance of equijoins in a
variety of scenarios. Furthermore, this paper presents a low-level library that
helps users customise and define their own synthetic data sources for future
benchmarking purposes.

In their distinguished paper
\relalg{}~\cite{RelationalAlgebraByWayOfAdjunctions} it was noted that the
monadic structure of bulk types can help explain most of relational algebra.
Using this structure, the authors designed a new method to facilitate
the use of monad comprehensions in an efficient implementation of equijoins of
relational databases. This project presents an implementation of such a system
and an evaluation of the performance such query optimisations carry using the
bespoke tooling described above. 
\end{abstract}

\tableofcontents

% Main matter 
\chapter{Introduction} % 1-3 pages
\begin{comment}
It’s a good idea to *try* to write the introduction to your final report early on in the project. However, you will find it hard, as you won’t yet have a complete story and you won’t know what your main contributions are going to be. However, the exercise is useful as it will tell you what you *don’t* yet know and thus what questions your project should aim to answer. For the interim report this section should be a short, succinct, summary of the project’s main objectives. Some of this material may be re-usable in your final report, but the chances are that your final introduction will be quite different.  You are therefore advised to keep this part of the interim report short, focusing on the following questions: What is the problem, why is it interesting and what’s your main idea for solving it?  (DON'T use those three questions as subheadings however!  The answers should emerge from what you write.)
\end{comment}

Databases are absolutely vital to modern day society and contain domain specific knowledge about anything from specialised images of eyes \todo{add citation here} to the structure of crystals \cite{CambridgeStructuralDatabase}. Since its conception many different data models describing how to hold the data in databases have emerged, including the relational model \cite{RelationalModel} and the semi-structured model\cite{DatabaseSystems}.

In this project we concern ourselves with the relational model (to be introduced in \fref{sec:relationalmodel}) as a way of modelling the database. This rich model has many methods of expressing queries, especially relational algebra and relational calculus, both with their strengths and weaknesses \cite{RelationalCalculus,RelationalModel}. \todo{they are equivalent though?} However, the favoured specification of the authors of \cite{RelationalAlgebraByWayOfAdjunctions} seemed to be through list comprehensions, a beautiful feature that ``provide for a concise and expressive notation for writing list-processing code''. \cite{MonadComprehensions} They eventually propose using GHC's extended list comprehension syntax specifically designed to help bridge the already close relationship between relational calculus and list comprehensions \cite{GHCListComprehension,ComprehensiveComprehensions} in order to avoid the significant theoretical performance hit.

It is widely noted that \emph{joins}, an integral operation in relational algebra are associated with inefficient implementations \cite{JoinProcessing}. It is easy to see why when considering the most general joins. However in their paper \cite{RelationalAlgebraByWayOfAdjunctions} they concern themselves with a specialised join called an \emph{equijoin}. As described in \fref{sec:joins} an equijoin is a specialised \emph{theta-join} -- a way of combining two relations based off of an arbitrary condition depending on the attributes of both relations. An integral part to the calculation of a theta-join is calculating the Cartesian product (all possible combinations of tuples of both relationships, as described in \fref{sec:products}). The algorithm must then filter every single tuple individually to check for equality of attributes! It is clear that this is wasteful for such an specialised join. As a more practical example consider the SQL program:
\begin{lstlisting}
  SELECT *
  FROM R, S
  WHERE R.a = S.b
\end{lstlisting}
This could be naively converted into a list comprehension with the following:
\[
  \left[\,(r, s)\;|\;r \leftarrow R,\;s \leftarrow S,\;r.a = s.b\,\right]
\]
where $(r, s)$ is seen as a single tuple whose attributes \attribute{\relation{R}.a} and \attribute{\relation{R}.b} are merged. \todo{fix example so that you do not need to see things}

 With this naive list comprehension implementation we effectively convert \equijoin{R}{\attribute{a}}{S}{\attribute{b}} to \select{\attribute{a} = \attribute{b}}{\relation{R} \times \relation{S}}, generating a relation with $|R||S|$ tuples in the process then filtering through each.

 This can much be much more efficiently implemented by viewing databases as indexed tables. We can index each relation by its associated attribute in the equijoin and merge the results -- localising the data required to in a cartesian product. This approach admits a linear time equijoin, if careful about comparison and projection functions. \cite{RelationalAlgebraByWayOfAdjunctions} With some mathematical tools explained in \fref{sec:gradedmonads} we can describe give these operations a monadic structure and therefore a comprehension syntax using the extended syntax discussed above.

 What this project adds to this story is a concrete demonstration of the improvement this solution offers.
 We will use the \emph{Haskell} and the list comprehensions and functions described above to implement a simple database querying software taking into account these key changes. Along with real world pragmatic data sources, benchmarking techniques found in \fref{sec:benchmarking} will be used to accurately measure and compare the efficiency difference between the two approaches -- with the new equijoin and without. This evidence would be very important to justify the use of these methods and the claims made in the paper \cite{RelationalAlgebraByWayOfAdjunctions}.
 With a concrete implementation, it could also provide insights into the downfall of remaining operations as well as those mentioned in the paper, potentially using profiling techniques to further analyse the performance bottlenecks, in order to inspire a theoretical approach at determining the issue as shown in \cite{RelationalAlgebraByWayOfAdjunctions} instead of an efficiency-driven optimisation approach.

\chapter{Background} % 10-20 pages
\begin{comment}
This should form the bulk of the interim report. You should consider that your objective here is to produce a near final version of the background section, as it will appear in your final report.  All of this material should be re-usable, so it is worth getting it right at this stage of the project.  The details of what to include can be found in the Project Report guidelines.

From the Project Report guidelines:
The background section of the report should set the project into context by relating it to existing published work which you read at the start of the project when your approach and methods were being considered. There are usually many ways of solving a given problem, and you shouldn't just pick one at random. Describe and evaluate as many alternative approaches as possible. The published work may be in the form of research papers found in the academic literature, articles, text books, technical manuals, or even existing software or hardware of which you have had hands-on experience. You must acknowledge the sources of your inspiration. You are expected to have seen and thought about other people's ideas; your contribution will be putting them into practice in some other context. However, avoid plagiarism: if you take another person's work as your own and do not cite your sources of information/inspiration you are being dishonest. When referring to other pieces of work, cite the sources where they are referred to or used, rather than just listing them at the end. Accidental plagiarism or not knowing how to cite and reference is not a valid reason for plagiarism. Make sure you read and digest the Department's plagiarism document .

In writing the Background chapter you must demonstrate your ability to *analyse*, *synthesise* and *apply critical* judgement. Analysis is shown by explaining how the proposed solution operates in your own words as well as its benefits and consequences. Synthesis is shown through the organisation of your Related Work section and through identifying and generalising common aspects across different solutions. Critical judgement is shown by discussing the limitations of the solutions proposed both in terms of their disadvantages and limits of applicability.

Typically you can look for Background work using different search engines including:
* Google Scholar
* IEEExplore
* ACM Digital Library
* Citeseer
* Science Direct

**Note 1:** Often the terms *Background, Literature Review, Related Work* and *State of the Art* are used interchangeably.
**Note 2**: Keyword search is wonderful, but you need the right *Keywords*.
**Note 3:** IEEExplore, ACM Digital Library and Science Direct may require you to be on the College network to download the PDF versions of papers. If at home, use VPN.
\end{comment}

\section{Crash course}
\subsection{Category theory}
\subsubsection{Categories}
\subsubsection{Functors}
\subsubsection{Natural transformations}
\subsubsection{Adjunctions}
\todo{Add more information between here}
\subsubsection{Graded Monads}

\subsection{Relational algebra}
\todo{Add citation to the database book}
\subsubsection{Projections}
\subsubsection{Selections}
\subsubsection{Products}
\subsubsection{Joins}

\subsection{Evolution of database representation}
\subsubsection{Bags}
\todo{Understand and distinguish between bags being the bulk type}
\paragraph{Characteristics of a database}We expect our database approximation to not be ordered and admit multiplicities and a finite bag of values is one of the simplest constructions that does so. Like a finite set, a bag contains a collection of unordered values. However, unlike a set, bags can contain duplicate elements. \cite{RelationalAlgebraByWayOfAdjunctions} This multiplicty is key for processing non-idempotent aggregations. For instance, if summing up the ages of a database of people, without admitting multiplicty we would only sum each unique age once.
\subparagraph{Generalisation}Furthermore, going forward we generalise to bags of any types instead of the classical ``bags of records''. This also allows us to deal with intermediate tables that contain non-record values.
\todo{Add the mathematical parts about bags}
\todo{When reading about fininte maps it says that it's better for databases as non finite maps cannot be aggreagated, can non finite bags be aggreagates - is this why they are finite?}

\subsubsection{Indexed tables}
We want to move towards an indexed representation of our table in order to equijoin by indexing. \todo{Understand if this is right and equijoin by indexing}. So in this section we introduce the mathematical concepts required to define such an implementation.
\theoremstyle{definition}\newtheorem*{psetdef}{Pointed set}
\theoremstyle{definition}\newtheorem*{ppfuncdef}{Point-preserving function}
\theoremstyle{definition}\newtheorem*{mapdef}{Map}
\theoremstyle{definition}\newtheorem*{finitemapdef}{Finite map}
\begin{psetdef}\label{def:pset}
  A pointed set $\pset{A}{a}$ is a set $A$ with a distinguished element $a \in A$.
\end{psetdef}
We commonly refer to the distinguished element of a set $A$ as $null_A$ or, when not ambiguous to do so, $null$.
\begin{ppfuncdef}\label{def:ppfunc}
  Given two pointed sets $\pset{A}{null_A}$ and $\pset{B}{null_B}$, a total function $f: A \rightarrow B$ is point-preserving if $f(null_A) = null_B$.
\end{ppfuncdef}
\todo{See if there is a way to have better spacing in brackets}

\todo{Add in more mathematical detail for point preserving functions}

We now have the mathematical tools required to define a map. In its finite form a map is widely known in computer science by many other names such as a dictionary, association lists or key-value maps.

Let $\keyset$ be a set and $\valset$ a pointed set. To those already familiar with maps, it may help to think of $\keyset$ as keys and $\valset$ as values.
\begin{mapdef}
  A map of type $\map{\keyset}{\valset}$ is a total function from K to V.
\end{mapdef}
\begin{finitemapdef}
  A finite map of type \finitemap{\keyset}{\valset} is a map where only a finite number of keys are mapped to $null_\valset$ (where $null_\valset$ is the distinguished element of \valset). 
\end{finitemapdef}
The advantage of using a finite map in a database is to allow aggregation.
\todo{Understand why only semimonoidal}
\todo{Introduced indexed tables}

\section{Literature Review}
\subsection{Benchmarking databases}
\subsection{Expanding the use of adjunctions}
\chapter{Project Plan} % 1-2 pages
\begin{comment}
You should explain what needs to be done in order to complete the project and roughly what you expect the timetable to be. Don’t forget to include the project write-up (the final report), as this is a major part of the exercise. It’s important to identify key milestones and also fall-back positions, in case you run out of time.  You should also identify what extensions could be added if time permits.  The plan should be complete and should include those parts that you have already addressed (make it clear how far you have progressed at the time of writing).  This material will *not* appear in the final report.
\end{comment}

The project has a clear starting point and a wide range of potential branches to eventually explore. This allows the project to venture down many different specific roots, ranging in application focused to theoretically based.
We outline the stages the project can take below.

\section{Initial research and literature review}\label{sec:researchstage}
As the project could eventually take a turn into the theoretical side of the paper \cite{RelationalAlgebraByWayOfAdjunctions} this part is crucial. The time frame for this will be ongoing and has started, however the bulk should be finished by the end of January. An intuition as well as knowledge of the applications of category theory in the wider range of computer science would be vital in order to spot potential improvements in any results found while implementing and benchmarking the findings of the base paper \cite{RelationalAlgebraByWayOfAdjunctions} .

I do expect, however, that there would be a period of time after finishing \fref{sec:implementation} to return to this step in more detail than normal in order to fully solidify the direction of the project.

Of course research is also important to create a fair and rigorous benchmarking system. This style of research should be concluded in the coming days in order to ensure that any implementation is built to play nicely with the tools to be used (and potentially initially specifically for the datasets chosen).
\section{Implementation of the differing equijoin algorithms}\label{sec:implementation}
After understanding the basics of the theoretical implications as well as how category theory underpins Haskell as whole, work can be started on the first tangible step. That is to implement a mini database system that can query a given dataset, both using the inefficient and efficient equijoins.
\subsection{Implementation}
Initially the implementation will be very tailored to the data sources chosen. This will be done to avoid the difficulties of generalising types and creating a coherent system under Haskell's strict rules. It is not too difficult or even out of the scope of the project, but I feel as though focusing on the results of benchmarking will be more efficient in earlier stages. I hope to get this done by the end of February.
\subsection{Evaluation and analysis}
The analysis should begin as soon as the implementation has been completed. At the very least, benchmarking seems to be quite heuristic and experience helps a lot, thus some time should be used to make sure that the results are as accurate as possible. At this stage, if experience has suggested that changing the implementation would not be too much effort, I might also consider changing the data source the tests were based on in order to deal with more interesting or accurate results if any issues seem to emerge. This process should hopefully take a couple weeks.

Evaluation, however, should be done as an ongoing process during the project. A plan to evaluate will be presented in \fref{sec:correctnessevaluation}.

\section{Theoretical analysis of the remaining relational algebra}\label{sec:theoreticalanalysis}
After the implementation of the findings of the paper, I hope to have gained a much deeper understanding of the practicalities of database querying. This along with the mathematical intuition I hope to have gained during the stage described \fref{sec:researchstage} would inspire thought to which other operations currently do not adhere to a monadic structure or efficient implementations. Of course this would be greatly aided by research in the large field of database queries, potentially choosing to describe already known optimisations using this new model and adjunctions instead. This section of the project could last until new work is no longer being found.
\section{Reporting the project}
% Report due Monday 19 June 2022
Reporting the project is a necessity. My project is unique in that the first section is itself an analysis and therefore, if done correctly, will almost immediately be written in a final report style while it is being conducted. Later parts of the project might need their own time set apart to write about as mathematics is rarely linear and the report should tell a compelling story.

My plan is to continue to develop the tools to efficiently write my report alongside the implementation of the database, and begin to draft and validate the final parts of the report a month before the deadline. Useful tools I have already integrated are both manual and commit triggered checks on the spelling and code style of the reports. I have also worked on automatic releases and management of different report versions, early in December. A few of the systems are helpful but not at their potential yet and so I might put some time to think about possible extensions, for instance custom dictionaries for the spell checker so it does not flag up domain specific terminology as errors. I definitely need to migrate the report structure I have to lhs2\TeX{} so that when I have code snippets to report it can be done with ease. This should be finished by February, in time to develop the database in the report if that is the workflow I find most beneficial.
\section{Possible extensions}
\subsection{Progress in optimising other aspects of relational algebra}
Given a successful stage in \fref{sec:theoreticalanalysis} this might give the project a lot of scope to grow as an extension over the final months. I could choose between using the category theory learned to write proofs of the theoretical advantage of any novel approach I think of, or just a novel description of what is already there. Alternatively I could extend my implementation to include the modified operations and do another analysis of their practical performance change.
\subsection{A full implementation}
If I choose to take this project down a more practical root, a full general implementation of the database system implemented could be made as an extension. This would allow a specific schema and interface with the data stored in Haskell's data structures itself, differing from database abstraction layers that communicate to SQL servers currently used in production.\cite{HDBC} Although this extension would not be production ready, I hope that it would be demonstrative on whether list comprehensions have a place among relational calculus, both by its viability to be implemented but also developer satisfaction using them.
\subsection{Applications of related fields to relational calculus}
Category theory is extremely general and has countless applications in computer science, from describing lambda calculi to describing type systems.\cite{BasicCategoryTheoryforCS} A particularly interesting paper about modelling concurrency \cite{CategoriesForModellingConcurrency} caught my eye as a topic to research into in the future, given how vital concurrency is to database management systems \cite{DatabaseSystems}.

Research could also be done into the viability of modelling mutating behaviours via referential calculus, that either change schema or components.

\chapter{Evaluation plan} % 1-2 pages
\begin{comment}
Project evaluation is very important, so it's important to think now about how you plan to measure success. For example, what functionality do you need to demonstrate?  What experiments to you need to undertake and what outcome(s) would constitute success?  What benchmarks should you use? How has your project extended the state of the art?  How do you measure qualitative aspects, such as ease of use?  These are the sort of questions that your project evaluation should address; this section should outline your plan.
\end{comment}
\section{During implementation}\label{sec:evaluationduringimplementation}
\chapter{Ethical issues} % 1-2 pages
\begin{comment}
What are the wider ethical, legal, professional and societal issues surrounding your project and the accompanying research? You should use the ethics checklist as the basis for this discussion. 
\end{comment}
The project itself (in its early stages) does not have a novel ideal as a useable output and thus is very neutral in its ethical impact. Instead it has the ability to recommend or withhold recommendation of an earlier approach to database optimisations via implementation and benchmarking results, and in its final stages may recommend further optimisations and so the more interesting ethical discussion is around the ethics of database optimisations.

Databases are a necessity for most modern applications. Furthermore, database optimisations are especially important for large, scaling products or services. This means that improving this area of computer science has a very indiscriminate effect on the type of software products that can make use of it. Of course, just as any private company could make use of the newer more efficient techniques, any military hosting a database could too. It is also worth noting the implementation is in a functional programming language, a rising paradigm \todo{cite} typically associated with reliability \todo{add another quality} \todo{Add citation for this} which aligns very well with the real-time and safety critical nature of military applications. \todo{cite that these are important for military applications}

In fact, it could be argued that any findings could benefit society at large in an environmental standpoint. The heart of the project is in database optimisations, which in turn reduce the amount of time required to process queries. \todo{research this} One could extrapolate, that without much of an increased processing complexity and resource consumption, this would lead to a reduction in the energy required for these database queries. Given how commonplace database queries are \todo{find stats} this could have a non-trivial impact.

It is worth noting that the early stages of the project involves implementing the ideas created in \cite{RelationalAlgebraByWayOfAdjunctions}. This causes no legal issues as there is no protection on an idea, lacking the tangibility (and ability to be patented) \todo{See exact laws}. On an ethical level, despite implementations, improvements and building on work is commonplace, encouraged and vital to the academic community, one of the main contributors of the paper is my supervisor and I interpret this as implicit consent. With good referencing practice, plagiarism would also be completely avoided.

There is an interesting discussion to be had on the ethical impact of the benchmarking dataset to be used in the early stages of the project. In order to evaluate the suggestions in \cite{RelationalAlgebraByWayOfAdjunctions} a database system will need to be implemented, but more consequentially one or more data sources will need to be used to evaluate the queries. I do not think that creating a novel data source will be very useful or worthwhile so I will need to acquire at least one somehow, raising at least two issues.
\subparagraph*{Data protection} The immediate issue that springs to mind would be whether the data source contains personal data and if so the subsequent data protection methods required to fulfil my legal and ethical obligation. To avoid this issue for primary data sources, I will not collect data when creating a benchmarking database, instead opting to randomly generate records. For data sources I acquire, they will all be in public domain already, and avoid sensitive personal data. This should minimise the risk to individuals affected. The data protection laws consider "organisation" and "structuring" as processing personal data, hence why this is relevant. The databases acquired might need to be reorganised into the data structured explained in this report \todo{cite}. \todo{See if this is enough} \todo{Cite data protection laws}
\subparagraph*{Database rights} Another consideration when outsourcing database collections, is the extent to which those who have created the database are protected. The creation of a database, although not creative, has been recognised as taking significant work and so those that create the database have a form of copyright protection on their work. \todo{Make sure that this is accurate, cite if so} This risk will be mitigated by ensuring that all databases used are in the public domain and the correct licenses are acquired.

\todo{Potentially comment on the distribution of my implementation}

\section{Extensions}
Furthermore, any domains investigated in this project have their own set of ethical considerations which reasonably cannot be discussed now (as they are largely unknown). Care should be taken on thinking of the ethical impact of the work that is being taken when branching out to other scenarios that use bulk-types. For instance, the suggested family of languages -- logic programming -- has close ties with artificial intelligence which itself is an ethically complex subject. \todo{Cite AI ethics} I think it is fair to say, however, that the work done will not bias any field in a particular way (ethical or not), but potentially contribute to more efficient implementations or different perspective of ideas within them.

% Back matter
\bibliographystyle{unsrtnat}
% DoC uses the Vancouver Referencing Format.
\bibliography{report/bibs/interim}

\end{document}
\chapter{PROJECT X}
\chapter{Evaluation}
\chapter{Conclusion}

% Back matter
\bibliographystyle{unsrtnat}
% DoC uses the Vancouver Referencing Format.
\bibliography{bibs/combined}

\end{document}
